\documentclass[11pt,a4paper]{article}
\usepackage[margin=1in]{geometry}
\usepackage{kotex}
\usepackage{amsmath,amssymb}
\usepackage{booktabs}
\usepackage{hyperref}
\usepackage{enumitem}
\usepackage{xcolor}
\usepackage{listings}

\lstset{
  language=Python,
  basicstyle=\ttfamily\scriptsize,
  commentstyle=\color{green!40!black},
  keywordstyle=\color{blue!70!black},
  stringstyle=\color{orange!60!black},
  showstringspaces=false,
  columns=fullflexible,
  keepspaces=true,
  frame=single,
  breaklines=true
}

\title{Hands-On Graph Neural Networks Using Python\\Chapter 7 상세 정리}
\author{A7610 그래프신경망특론}
\date{2026년 4월 15일}

\begin{document}
\maketitle
\tableofcontents
\newpage

\section{Chapter 7 개요}

Chapter 7의 핵심은 \textbf{Graph Attention Networks (GAT)}이다. 기존 메시지 패싱 기반 GNN과 달리, GAT는 이웃마다 다른 중요도 계수를 부여해 집계한다.

\begin{quote}
"같은 이웃 집계라도 모든 이웃을 똑같이 보지 않고, 더 중요한 이웃에 더 큰 가중치를 준다."
\end{quote}

실습 노트북 \texttt{practice/hands\_on\_gnn\_examples/Chapter07/chapter7.ipynb} 흐름은 다음과 같다.
\begin{enumerate}[leftmargin=*]
  \item toy 예제로 attention 계산 절차를 수치적으로 확인
  \item PyTorch Geometric으로 GATv2 모델 구현
  \item Cora/CiteSeer 데이터셋에서 분류 성능 비교
  \item 노드 degree별 정확도 분석으로 성능 편차 해석
\end{enumerate}

\subsection*{노트북의 의도}
\texttt{chapter7\_annotated.ipynb}는 원본 실습 코드에 설명을 삽입한 버전이다.
\begin{itemize}[leftmargin=*]
  \item 코드 셀은 그대로 실행 가능하게 유지한다.
  \item 각 코드 블록 앞에 ``왜 이 계산이 필요한가''를 설명한다.
  \item 교재의 핵심 그림 6개를 함께 배치해 수식--코드--도식을 연결한다.
\end{itemize}

\subsection*{먼저 정리할 오해: GAN이 아니라 GAT}
발표 초반에는 용어 혼동을 먼저 끊는 편이 좋다.
\begin{itemize}[leftmargin=*]
  \item \textbf{GAN}: 생성자와 판별자가 경쟁하는 생성 모델
  \item \textbf{GAT}: 그래프 이웃의 중요도를 attention으로 학습하는 메시지 패싱 모델
\end{itemize}
즉, 이 장의 관심사는 ``가짜 샘플 생성''이 아니라 ``어떤 이웃의 메시지를 더 크게 들을 것인가''이다.

\section{GAT의 수학적 직관}

노드 $i$의 새로운 표현은 이웃 메시지의 가중합으로 표현된다.
\[
\mathbf{h}_i' = \sigma\left(\sum_{j\in\mathcal{N}(i)} \alpha_{ij}\mathbf{W}\mathbf{h}_j\right)
\]

여기서 attention 계수는 다음과 같이 계산된다.
\[
e_{ij} = \text{LeakyReLU}(\mathbf{a}^{\top}[\mathbf{W}\mathbf{h}_i\,\|\,\mathbf{W}\mathbf{h}_j]),
\quad
\alpha_{ij}=\frac{\exp(e_{ij})}{\sum_{k\in\mathcal{N}(i)}\exp(e_{ik})}
\]

핵심은 \textbf{중심 노드별 softmax 정규화}다. 즉, 노드 $i$의 이웃들 내부에서 상대적으로 중요한 이웃이 높은 계수를 받는다.

\section{노트북 1부: toy 예제 해설}

노트북은 4개 노드 인접행렬 $A$를 직접 정의하고, 랜덤 초기화된 특징행렬 $X$, 투영행렬 $W$, attention 벡터 $W_{att}$를 사용해 계수 계산을 시뮬레이션한다.

\begin{lstlisting}
connections = np.where(A > 0)

a = W_att @ np.concatenate([
    (X @ W.T)[connections[0]],
    (X @ W.T)[connections[1]]
], axis=1).T

e = leaky_relu(a)
E = np.zeros(A.shape)
E[connections[0], connections[1]] = e[0]
W_alpha = softmax2D(E, 1)
\end{lstlisting}

이 구현은 논문 수식을 코드로 어떻게 번역하는지 보여준다.
\begin{itemize}[leftmargin=*]
  \item 연결된 노드쌍만 점수를 계산한다.
  \item LeakyReLU 후 각 중심 노드에서 softmax를 적용한다.
  \item 최종적으로 가중치 행렬 $W\_\alpha$를 얻는다.
\end{itemize}

\subsection*{toy example를 읽는 순서}
노트북의 설명을 발표용 언어로 옮기면 다음 순서가 가장 명확하다.
\begin{enumerate}[leftmargin=*]
  \item \textbf{입력 준비}: $A$는 연결 구조, $X$는 노드 특징, $W$는 특징 변환, $W_{att}$는 attention 점수 생성 파라미터다.
  \item \textbf{연결된 노드쌍 추출}: \texttt{connections = np.where(A > 0)}로 edge가 있는 쌍 $(i,j)$만 남긴다.
  \item \textbf{원시 점수 계산}: \texttt{np.concatenate(...)}가 $[\mathbf{W}\mathbf{h}_i \, \| \, \mathbf{W}\mathbf{h}_j]$에 대응한다.
  \item \textbf{정규화}: \texttt{softmax2D(E, 1)}에서 축 1은 ``한 중심 노드의 이웃 집합''을 뜻한다.
  \item \textbf{가중합}: attention이 큰 이웃일수록 최종 표현 $H$에 더 크게 반영된다.
\end{enumerate}

이 과정을 손으로 따라가면 뒤에서 \texttt{GATv2Conv}가 수행하는 일이 블랙박스가 아니라는 점이 중요하다.

\subsection*{교재 그림과 코드 연결}
교재 그림 7.1--7.2는 노트북 코드와 거의 일대일로 대응된다.
\begin{itemize}[leftmargin=*]
  \item \texttt{connections = np.where(A > 0)}: 실제 연결된 노드쌍만 추출
  \item \texttt{a = W\_att @ concat(...)}: edge별 raw score 생성
  \item \texttt{e = leaky\_relu(a)}: 비선형 점수화
  \item \texttt{W\_alpha = softmax2D(E, 1)}: 이웃 축 기준 정규화
\end{itemize}

\section{노트북 2부: Cora/CiteSeer 실험}

PyG의 \texttt{Planetoid} 데이터셋을 사용해 Cora와 CiteSeer를 불러온다.

\begin{lstlisting}
dataset = Planetoid(root=".", name="Cora")
data = dataset[0]
\end{lstlisting}

모델은 2층 GATv2Conv로 구성된다.

\begin{lstlisting}
class GAT(torch.nn.Module):
    def __init__(self, dim_in, dim_h, dim_out, heads=8):
        super().__init__()
        self.gat1 = GATv2Conv(dim_in, dim_h, heads=heads)
        self.gat2 = GATv2Conv(dim_h*heads, dim_out, heads=1)
\end{lstlisting}

Forward에서는 dropout-eLU-dropout-GAT의 전형적 구조를 사용한다.

\begin{lstlisting}
h = F.dropout(x, p=0.6, training=self.training)
h = self.gat1(h, edge_index)
h = F.elu(h)
h = F.dropout(h, p=0.6, training=self.training)
h = self.gat2(h, edge_index)
\end{lstlisting}

\subsection*{PyG 구현에서 짚고 넘어갈 포인트}
\begin{itemize}[leftmargin=*]
  \item \textbf{첫 층 multi-head}: 여러 head가 서로 다른 방식으로 이웃 중요도를 보게 해 표현력을 높인다.
  \item \textbf{두 번의 dropout}: 입력 특징과 은닉 표현 모두에 regularization을 걸어 Planetoid 계열 데이터셋에서 과적합을 줄인다.
  \item \textbf{학습 마스크}: 손실은 train mask에 있는 노드에 대해서만 계산되므로, 전체 그래프를 보더라도 supervised signal은 제한적이다.
\end{itemize}

노트북 설명을 발표용으로 정리하면, ``attention 메커니즘 자체''보다도 \textbf{작은 라벨 수와 작은 그래프에서 일반화가 되도록 어떤 학습 장치를 같이 쓰는가}가 실험 성능을 좌우한다.

\subsection*{실험 설정}
\begin{itemize}[leftmargin=*]
  \item 데이터셋: Cora, CiteSeer
  \item optimizer: Adam, learning rate $0.01$, weight decay $0.01$
  \item epoch: 100
  \item dropout: 0.6
  \item attention heads: 첫 층 8개, 출력층 1개
\end{itemize}

\section{결과 해석}

노트북 출력 기준 테스트 정확도는 다음과 같다.
\begin{itemize}[leftmargin=*]
  \item Cora: \textbf{81.10\%}
  \item CiteSeer: \textbf{68.10\%}
\end{itemize}

같은 모델이라도 데이터셋 특성에 따라 성능 차이가 크다. 이 차이는 다음 요인의 조합으로 설명할 수 있다.
\begin{itemize}[leftmargin=*]
  \item 그래프 구조 밀도와 커뮤니티 분리도
  \item 특징 벡터의 정보량
  \item train/val/test split의 난이도 차이
\end{itemize}

\subsection*{결과를 읽는 법}
노트북의 해설은 평균 정확도 숫자 하나만 읽지 말라고 강조한다.
\begin{itemize}[leftmargin=*]
  \item Cora와 CiteSeer의 성능 차이는 모델 차이만이 아니라 데이터셋 성질 차이를 반영한다.
  \item 따라서 ``GAT가 항상 더 좋다''가 아니라 ``어떤 구조에서 attention이 더 잘 작동하는가''를 묻는 것이 맞다.
  \item 교재 그림 7.5--7.6은 성능 숫자보다 해석 포인트와 한계를 읽는 데 초점을 둔다.
\end{itemize}

\section{Degree별 성능 분석}

노트북 후반부는 노드 degree 구간별 정확도를 계산한다.

\begin{lstlisting}
degrees = degree(data.edge_index[0]).numpy()

for i in range(0, 6):
    mask = np.where(degrees == i)[0]
    accuracies.append(accuracy(out.argmax(dim=1)[mask], data.y[mask]))

mask = np.where(degrees > 5)[0]
accuracies.append(accuracy(out.argmax(dim=1)[mask], data.y[mask]))
\end{lstlisting}

이 분석은 평균 정확도 한 줄로 보이지 않는 현상을 보여준다.
\begin{itemize}[leftmargin=*]
  \item 저차수 노드: 이웃 정보가 부족해 예측이 불안정할 수 있다.
  \item 고차수 노드: 정보가 풍부하지만 과도한 혼합으로 잡음이 들어올 수도 있다.
\end{itemize}

\subsection*{왜 이 분석이 중요한가}
전체 정확도만 보면 다음 문제가 가려질 수 있다.
\begin{itemize}[leftmargin=*]
  \item 저차수 노드에서만 성능이 무너지는지
  \item 고차수 노드에서 과도한 혼합이 생기는지
  \item 특정 구조 구간이 모델의 취약점인지
\end{itemize}
즉, degree별 성능은 ``모델이 어디에서 실패하는가''를 구조적으로 진단하는 도구다.

\section{발제 포인트}

발표에서는 아래 질문을 던지면 토론이 좋아진다.
\begin{enumerate}[leftmargin=*]
  \item attention 계수가 항상 해석 가능한 중요도인가?
  \item heads 수를 늘리면 언제 성능이 좋아지고 언제 과적합되는가?
  \item heterophily 그래프에서도 동일한 attention 설계가 유효한가?
\end{enumerate}

\subsection*{노트북에서 바로 이어갈 추가 실습}
\begin{enumerate}[leftmargin=*]
  \item \texttt{heads=4,8,16}으로 바꿔 성능과 분산 비교
  \item \texttt{dropout=0.3,0.6,0.8} 비교
  \item 같은 split에서 \texttt{GATv2Conv}와 \texttt{GCNConv} 비교
  \item 오분류 노드의 degree와 이웃 라벨 분포 분석
\end{enumerate}

\section{Chapter 7 핵심 요약}

\begin{itemize}[leftmargin=*]
  \item GAT는 메시지 패싱에 학습형 이웃 가중치(attention)를 도입한다.
  \item toy 예제는 수식과 구현의 대응관계를 이해하는 데 유용하다.
  \item PyG 구현에서는 multi-head, dropout, train mask가 함께 작동해 일반화 성능을 만든다.
  \item Cora/CiteSeer 실험은 데이터셋 성질에 따른 성능 차이를 보여준다.
  \item degree별 성능 분석은 모델의 구조적 편향을 진단하는 좋은 도구다.
\end{itemize}

\end{document}
